\documentclass[../../../include/open-logic-section]{subfiles}

\begin{document}

\olfileid[fr]{his}{set}{cantorplane}
\olsection{Cantor, la droite et le plan}

Certaines circonstances entourant la démonstration du théorème de
Schr\"oder--Bernstein se rattachent à l'histoire évoquée dans la
\olref[his][set][pathology]{sec}. Rappelons qu'en 1877, Cantor
démontra qu'un carré contient exactement autant de points que
l'un de ses côtés. Nous allons présenter sa première tentative
de démonstration.

Soit $\unitline$ le segment unité, c'est-à-dire l'ensemble de
points $[0,1]$, et soit $\unitsquare$ le carré unité, c'est-à-dire
$\unitline \times \unitline$. Dans cette notation, Cantor démontra
que $\cardeq{\unitline}{\unitsquare}$. Il envoya à Dedekind une
note contenant, en substance, l'argument suivant.

\begin{thm}\ollabel{thm:cantorplane}
$\cardeq{\unitline}{\unitsquare}$
\end{thm}

\begin{proof}[Première partie de la démonstration]
Fixons $a, b \in \unitline$ et écrivons-les en notation binaire.
Pour tout nombre strictement inférieur à~$1$, choisissons le
développement qui n'est pas ultimement constitué uniquement de~$1$ ;
en particulier, choisissons le développement terminé par des~$0$
lorsqu'il existe. Écrivons $1$ sous la forme $0.111\ldots$.\footnote{
Précision de l'édition française : la source anglaise ne fixe pas
la convention de développement binaire. La convention donnée ici
définit l'application sur tout le segment $[0,1]$ et permet de
justifier son injectivité.}
Nous obtenons ainsi des suites infinies de~$0$ et de~$1$,
$a_1$, $a_2$, \dots, et $b_1$, $b_2$, \dots, telles que :
\begin{align*}
a &= 0.a_1a_2a_3a_4\dots\\
b &= 0.b_1b_2b_3b_4\dots
\intertext{Considérons l'application $f \colon \unitsquare \to \unitline$ définie par}
f(a, b) & = 0.a_1b_1a_2b_2a_3b_3a_4b_4\dots
\end{align*}
Le développement entrelacé n'est ultimement constitué uniquement
de~$1$ que si $a=b=1$, auquel cas $f(a,b)=1$. Les images sont donc
écrites selon la même convention. Par conséquent, $f$ est
!!a{injection} : si $f(a, b) = f(c,d)$, l'unicité de ces
développements donne $a_n = c_n$ et $b_n = d_n$ pour tout
$n \in \Nat$, d'où $a = c$ et $b = d$.
\end{proof}

Malheureusement, comme Dedekind le fit remarquer à Cantor, cela ne
répond pas à la question initiale. Avec la convention ci-dessus,
considérons $y=0.00010101010\ldots$.\footnote{
Correction mathématique de l'édition française : l'exemple de la
source, $0.\dot{1}\dot{0}$, est bien l'image de $(1,0)$ lorsque
$1$ est représenté par $0.111\ldots$, comme l'exige notre convention.
Nous utilisons donc un autre développement pour montrer le défaut
de surjectivité.}
Pour avoir $f(a,b)=y$, il faudrait que les chiffres binaires
de $a$ et de $b$ soient respectivement :
\begin{align*}
a&=0.000000\ldots=0,\\
b&=0.011111\ldots=\tfrac{1}{2}.
\end{align*}
Mais le développement choisi pour $\tfrac{1}{2}$ est
$0.100000\ldots$, et non $0.011111\ldots$. Ainsi,
$f(0,\tfrac{1}{2})=0.010000\ldots$, et aucun couple
$\tuple{a,b}$ n'a pour image~$y$. Le choix entre les deux
développements possibles est essentiel : pour que $f$ soit une
\emph{fonction}, il faut n'en utiliser qu'\emph{un} pour chaque
nombre.

Le problème signalé par Dedekind tient aux deux développements
possibles d'un même nombre dyadique : l'application $f$
de Cantor est !!a{injection}, mais \emph{pas} !!a{surjection}.
Cantor avait seulement montré que
$\cardle{\unitsquare}{\unitline}$, et \emph{non} que
$\cardeq{\unitsquare}{\unitline}$.

Cantor répondit presque immédiatement à Dedekind en suggérant,
en substance, de compléter l'argument ainsi :

\begin{proof}[Conclusion de la démonstration]
Nous avons montré que $\cardle{\unitsquare}{\unitline}$.
Il existe manifestement !!a{injection} de $\unitline$ dans
$\unitsquare$ : il suffit de placer le segment le long d'un côté
du carré. Nous avons donc $\cardle{\unitline}{\unitsquare}$ et
$\cardle{\unitsquare}{\unitline}$. Par le théorème de
Schr\"{o}der--Bernstein
(\olref[sfr][siz][sb]{thm:schroder-bernstein}),
$\cardeq{\unitline}{\unitsquare}$.
\end{proof}

Bien sûr, Cantor ne pouvait pas formuler la dernière étape en ces
termes, puisque le théorème de Schr\"{o}der--Bernstein n'avait
pas encore été démontré. Cantor l'énonça ensuite comme une
conjecture générale, mais une démonstration satisfaisante ne fut
obtenue qu'en 1897. Plus tard en 1877, il donna donc une autre
démonstration du \olref{thm:cantorplane}, qui ne passait pas par
Schr\"{o}der--Bernstein.

\end{document}
